\section{Pending Coauthor Contributions and Joint Integration}\label{app:team-pending}
This working appendix reserves the missing methods and evidence for the shared project. The earlier preprocessing PDF is a provisional contribution, not a verified final supplement. Its numbers have not been copied into the joint results. No learning-curve outcomes are available for integration. TODOs are deliberate draft markers and must be resolved before final submission.

\subsection{RQ1: Preprocessing and Detection Recovery}
\drafttodo{PREP}{Provide the final text, source code/version and result artifacts. Confirm the baseline manifest, detector model/version, detection/presence thresholds and rejection categories. Explain which initially rejected images are eligible for retry.}
\drafttodo{PREP}{Define dilation and the structuring element, colour handling, resolution, iterations, erosion and Canny-overlay settings. Keep the proposed $3\times3$ recipe provisional until the author confirms it.}
\drafttodo{PREP}{Report original misses, recovered images per variant, unique union and remaining rejections, including per-letter counts. Individual variant sets may overlap: percentages cannot simply be added. State the denominator and source IDs for every result.}
\drafttodo{PREP}{Supply a recovery figure and visual-review method/results, including the review sample, criteria and assessor procedure. Technical acceptance is not proof of anatomically correct landmarks. If no visual audit is completed, retain that limitation explicitly.}

\subsection{RQ2: Amount of Training Data}
\drafttodo{SIZE}{Provide the exact source population, fitting sizes and counts per letter, accepted-detection counts, sampling unit and policy. State whether this uses only original accepted images or any recovered images; do not change that population silently across the curve.}
\drafttodo{SIZE}{Provide fixed held-out source IDs/folds, subset IDs, seeds and repeat counts. Draw fitting subsets within the training partition only, preferably nested within each repeat. Keep class balance controlled and explain any exclusions.}
\drafttodo{SIZE}{Specify detector, feature representation, classifier settings and training-only preprocessing at every size. Define precision averaging, F1 and uncertainty summaries. Repeated source subsets do not create independent test populations.}
\drafttodo{SIZE}{Supply the learning-curve plot and underlying values for macro-precision and macro-F1, repeat/fold variation and per-letter changes. Report plateaus or decreases if observed; a monotonic improvement is not assumed.}

\subsection{Shared Comparison and Handoff}
\drafttodo{TEAM}{Confirm author names, student numbers, ownership and the final wording of RQ1--RQ3. Add relevant preprocessing/learning-curve references and verify that cited work supports the claims.}
\drafttodo{TEAM}{Reconcile dataset versions, manifests, detector settings, source splits and metric populations. The historical feature cache does not retain detector thresholds; do not describe separately obtained landmarks as identical without checking provenance.}
\drafttodo{TEAM}{Use numbered tables/figures for findings and appendix sections for detail. Put the runnable code and versioned CSV/JSON evidence in the supplement with stable paths and checksums; the live application is a demonstration, not the numerical evidence.}
\drafttodo{TEAM}{Write the joint abstract, discussion and conclusion after the missing results arrive. Distinguish recovery, visual plausibility, conditional classification and all-attempt performance. Combining recovered data, larger fitting sets and different features would require a separate matched experiment; independently measured gains cannot be added.}

The editable insertion points are \texttt{team\_methods\_pending.tex} and \texttt{team\_results\_pending.tex}. \texttt{INTEGRATION\_TODO.md} maps the handoff to the common sections. The final main text remains limited to five pages excluding references, with detailed evidence in the appendix and supplement.
